\documentclass[aspectratio=169, 14pt]{beamer}
\usepackage{stampfly_slides}

\lesson{4}{IMU センサ}{IMU Sensor}

\begin{document}

% ------------------------------------------------------------------
\begin{frame}
  \titlepage
\end{frame}

% ------------------------------------------------------------------
\begin{frame}{今日のゴール / Today's Goal}
  \begin{block}{目標}
    IMU（ジャイロ+加速度）データを読み取り、テレメトリで可視化する
  \end{block}

  \vspace{1em}

  \begin{itemize}
    \item NED 座標系（北-東-下）の理解
    \item ジャイロスコープで角速度を取得
    \item 加速度センサで並進加速度を取得
    \item WiFi テレメトリでデータ取得・可視化
  \end{itemize}
\end{frame}

% ------------------------------------------------------------------
\begin{frame}{NED 座標系 / NED Coordinate System}
  \vspace{-0.5em}
  \begin{center}
    \resizebox{0.65\textwidth}{!}{\documentclass[border=5pt]{standalone}
\usepackage{tikz}
\usetikzlibrary{arrows.meta, calc, decorations.markings}

\begin{document}

\definecolor{axisX}{RGB}{220, 50, 30}   % Red - X forward (North)
\definecolor{axisY}{RGB}{40, 160, 40}   % Green - Y right (East)
\definecolor{axisZ}{RGB}{0, 120, 200}   % Blue - Z down
\definecolor{bodycolor}{RGB}{80, 80, 80}

\begin{tikzpicture}[
    axis/.style={-{Stealth[length=5mm, width=3mm]}, line width=2.5pt},
    rotarrow/.style={
      -{Stealth[length=3.5mm, width=2.5mm]},
      line width=1.8pt,
    },
  ]

  % Isometric unit vectors
  \pgfmathsetmacro{\Xx}{-0.71}
  \pgfmathsetmacro{\Xy}{0.50}
  \pgfmathsetmacro{\Yx}{1.0}
  \pgfmathsetmacro{\Yy}{0.0}
  \pgfmathsetmacro{\Zx}{0.0}
  \pgfmathsetmacro{\Zy}{-1.0}

  % Body - isometric parallelogram
  \fill[bodycolor!10]
    ({1.1*\Xx+1.1*\Yx},{1.1*\Xy+1.1*\Yy}) --
    ({1.1*\Xx-1.1*\Yx},{1.1*\Xy-1.1*\Yy}) --
    ({-1.1*\Xx-1.1*\Yx},{-1.1*\Xy-1.1*\Yy}) --
    ({-1.1*\Xx+1.1*\Yx},{-1.1*\Xy+1.1*\Yy}) -- cycle;
  \draw[bodycolor, line width=1.5pt]
    ({1.1*\Xx+1.1*\Yx},{1.1*\Xy+1.1*\Yy}) --
    ({1.1*\Xx-1.1*\Yx},{1.1*\Xy-1.1*\Yy}) --
    ({-1.1*\Xx-1.1*\Yx},{-1.1*\Xy-1.1*\Yy}) --
    ({-1.1*\Xx+1.1*\Yx},{-1.1*\Xy+1.1*\Yy}) -- cycle;

  % Front marker triangle
  \fill[axisX!50]
    ({1.1*\Xx-0.35},{1.1*\Xy}) --
    ({1.1*\Xx+0.35},{1.1*\Xy}) --
    ({1.1*\Xx},{1.1*\Xy+0.3}) -- cycle;

  % Axis length
  \pgfmathsetmacro{\AL}{4.5}

  % X axis (Forward / North)
  \draw[axis, axisX] (0,0) -- ({\AL*\Xx},{\AL*\Xy});
  \node[font=\Large\bfseries, axisX, above left]
    at ({\AL*\Xx},{\AL*\Xy}) {X (Forward)};

  % Y axis (Right / East)
  \draw[axis, axisY] (0,0) -- ({\AL*\Yx},{\AL*\Yy});
  \node[font=\Large\bfseries, axisY, above right]
    at ({\AL*\Yx},{\AL*\Yy}) {Y (Right)};

  % Z axis (Down)
  \draw[axis, axisZ] (0,0) -- ({\AL*\Zx},{\AL*\Zy});
  \node[font=\Large\bfseries, axisZ, right]
    at ({\AL*\Zx},{\AL*\Zy}) {Z (Down)};

  % === Rotation arrows (elliptical arcs wrapping around each axis) ===

  % Roll - rotation around X axis
  % Arc lies in the Y-Z plane, centered on X axis
  % Projected semi-axes: Y direction (1,0) and Z direction (0,-1)
  \pgfmathsetmacro{\Rcx}{2.6*\Xx}  % center on X axis
  \pgfmathsetmacro{\Rcy}{2.6*\Xy}
  \pgfmathsetmacro{\Rr}{0.9}        % arc radius
  \draw[rotarrow, axisX!70!black]
    ({\Rcx+\Rr*\Yx*cos(200)+\Rr*\Zx*sin(200)},
     {\Rcy+\Rr*\Yy*cos(200)+\Rr*\Zy*sin(200)})
    \foreach \a in {200,210,...,490} {
      -- ({\Rcx+\Rr*\Yx*cos(\a)+\Rr*\Zx*sin(\a)},
          {\Rcy+\Rr*\Yy*cos(\a)+\Rr*\Zy*sin(\a)})
    };
  \node[font=\normalsize\bfseries, axisX!70!black, above]
    at ({\Rcx},{\Rcy+\Rr+0.15}) {Roll};

  % Pitch - rotation around Y axis
  % Arc lies in the X-Z plane, centered on Y axis
  % Projected semi-axes: X direction (-0.71,0.5) and Z direction (0,-1)
  \pgfmathsetmacro{\Pcx}{2.6*\Yx}
  \pgfmathsetmacro{\Pcy}{2.6*\Yy}
  \pgfmathsetmacro{\Pr}{0.9}
  \draw[rotarrow, axisY!70!black]
    ({\Pcx+\Pr*\Xx*cos(200)+\Pr*\Zx*sin(200)},
     {\Pcy+\Pr*\Xy*cos(200)+\Pr*\Zy*sin(200)})
    \foreach \a in {200,210,...,490} {
      -- ({\Pcx+\Pr*\Xx*cos(\a)+\Pr*\Zx*sin(\a)},
          {\Pcy+\Pr*\Xy*cos(\a)+\Pr*\Zy*sin(\a)})
    };
  \node[font=\normalsize\bfseries, axisY!70!black, above right]
    at ({\Pcx+\Pr*0.3},{\Pcy+\Pr*0.7}) {Pitch};

  % Yaw - rotation around Z axis
  % Arc lies in the X-Y plane, centered on Z axis
  % Projected semi-axes: X direction (-0.71,0.5) and Y direction (1,0)
  \pgfmathsetmacro{\Wcx}{2.6*\Zx}
  \pgfmathsetmacro{\Wcy}{2.6*\Zy}
  \pgfmathsetmacro{\Wr}{0.9}
  \draw[rotarrow, axisZ!70!black]
    ({\Wcx+\Wr*\Xx*cos(160)+\Wr*\Yx*sin(160)},
     {\Wcy+\Wr*\Xy*cos(160)+\Wr*\Yy*sin(160)})
    \foreach \a in {160,170,...,450} {
      -- ({\Wcx+\Wr*\Xx*cos(\a)+\Wr*\Yx*sin(\a)},
          {\Wcy+\Wr*\Xy*cos(\a)+\Wr*\Yy*sin(\a)})
    };
  \node[font=\normalsize\bfseries, axisZ!70!black, left]
    at ({\Wcx-\Wr*0.8},{\Wcy+\Wr*0.3}) {Yaw};

\end{tikzpicture}
\end{document}
}
  \end{center}
  \vspace{-0.5em}
  \begin{block}{}
    \centering\small
    BMI270 IMU --- 6軸（3軸ジャイロ + 3軸加速度）400\,Hz
    \quad|\quad
    回転矢印 = ジャイロ正方向（右手の法則）
  \end{block}
\end{frame}

% ------------------------------------------------------------------
\begin{frame}{IMU API}
  \begin{table}
    \centering\small
    \begin{tabular}{lll}
      \hline
      \textbf{関数} & \textbf{説明} & \textbf{単位} \\
      \hline
      \api{gyro\_x/y/z()}  & 角速度 (Roll/Pitch/Yaw) & rad/s \\
      \api{accel\_x/y/z()} & 加速度 (X/Y/Z) & m/s\textsuperscript{2} \\
      \api{telemetry\_send(name, val)} & WiFi 送信 & --- \\
      \hline
    \end{tabular}
  \end{table}

  \vspace{0.5em}

  \begin{block}{静止時の値}
    ジャイロ $\approx 0$、加速度 Z $\approx 9.81$\,m/s\textsuperscript{2}（重力）
  \end{block}
\end{frame}

% ------------------------------------------------------------------
\begin{frame}[fragile]{実習: 6軸読み取り + テレメトリ / Hands-on}
  \begin{lstlisting}[style=sfcpp, title={student.cpp}]
#include "workshop_api.hpp"
static uint32_t tick = 0;
void setup() {
    ws::print("Lesson 4: IMU Sensor");
}
void loop_400Hz(float dt) {
    tick++;
    float gx = ws::gyro_x();
    float gy = ws::gyro_y();
    float gz = ws::gyro_z();
    // Send via WiFi telemetry
    ws::telemetry_send("gyro_x", gx);
    ws::telemetry_send("gyro_y", gy);
    ws::telemetry_send("gyro_z", gz);
    // Print every 200ms (80 ticks)
    if (tick % 80 == 0) {
        ws::print("gx=%.2f gy=%.2f gz=%.2f", gx, gy, gz);
    }
}
  \end{lstlisting}
\end{frame}

% ------------------------------------------------------------------
\begin{frame}{WiFi テレメトリ受信 / Receiving WiFi Telemetry}
  \begin{block}{手順}
    \begin{enumerate}
      \item シリアルモニタで SSID を確認（\code{StampFly\_XXXX}）
      \item PC の WiFi で StampFly に接続（IP: \code{192.168.4.1}）
      \item \code{sf log wifi -d 30} で30秒キャプチャ → CSV 自動保存
    \end{enumerate}
  \end{block}

  \vspace{0.5em}

  \begin{exampleblock}{出力例}
    \code{logs/stampfly\_wifi\_20260303T143000.csv} に保存されました
  \end{exampleblock}

  \vspace{0.3em}
  {\small
  \code{-o name.csv} でファイル名指定、\code{--no-save} で保存なし（統計のみ）
  }
\end{frame}

% ------------------------------------------------------------------
\begin{frame}{データの可視化 / Data Visualization}
  \begin{block}{保存した CSV を可視化}
    \code{sf log viz logs/stampfly\_wifi\_*.csv}
  \end{block}

  \vspace{0.5em}

  \begin{table}
    \centering\small
    \begin{tabular}{ll}
      \hline
      \textbf{モード} & \textbf{表示内容} \\
      \hline
      \code{--mode all}      & 全パネル（デフォルト） \\
      \code{--mode sensors}  & IMU + 高度センサ \\
      \code{--mode attitude} & 姿勢推定（Roll/Pitch/Yaw） \\
      \hline
    \end{tabular}
  \end{table}

  \vspace{0.5em}

  \begin{itemize}
    \item \code{--save plot.png} でファイル保存可能
    \item StampFly を手で揺らしながらキャプチャし、ジャイロの変化をグラフで確認しよう
  \end{itemize}
\end{frame}

% ------------------------------------------------------------------
\begin{frame}{チェックポイント / Checkpoint}
  \begin{alertblock}{確認事項}
    \begin{itemize}
      \item[$\square$] 手で傾けるとジャイロ値が変化する
      \item[$\square$] 静止時に accel\_z $\approx 9.81$ を確認
      \item[$\square$] \code{sf log wifi} でテレメトリデータを受信できる
      \item[$\square$] \code{sf log viz} でジャイロデータをグラフ表示できる
    \end{itemize}
  \end{alertblock}

  \vspace{1em}

  \begin{block}{次のレッスン}
    Lesson 5: レート P 制御 --- ジャイロフィードバックで初フライト
  \end{block}
\end{frame}

\end{document}
